\section{Проверка адекватности математических моделей}
\label{sec:model-adequacy-verification}

\subsection{Понятие адекватности}

\textbf{Адекватность математической модели} --- её способность с требуемой
точностью воспроизводить существенные для поставленной цели свойства реального
объекта в заданной области применимости. Поэтому модель нельзя называть
``адекватной вообще'': вывод всегда относится к определённым выходным величинам,
диапазону параметров, режимам и допустимой погрешности.

Пусть эксперимент даёт значения $y_i^{\rm exp}$, а модель ---
$y_i^{\rm mod}(\theta)$, где $\theta$ --- вектор параметров. Остатки
\[
    r_i=y_i^{\rm exp}-y_i^{\rm mod}(\theta)
\]
характеризуют расхождение расчёта и данных. Малая норма остатков необходима,
но сама по себе ещё не доказывает адекватность: необходимо учитывать погрешности
измерений, область проверки и то, использовались ли эти данные при настройке
параметров.

\subsection{Верификация, валидация и калибровка}

Следует различать три процедуры:
\begin{itemize}
    \item \textbf{верификация} --- правильно ли численно решена выбранная
    математическая задача;
    \item \textbf{валидация} --- достаточно ли хорошо сама математическая
    модель описывает реальный объект;
    \item \textbf{калибровка} --- определение параметров модели по данным.
\end{itemize}

Коротко:
\[
\boxed{\text{верификация: ``правильно ли мы решаем модель?''}},
\qquad
\boxed{\text{валидация: ``правильную ли модель мы решаем?''}}.
\]

Для верификации используют аналитические или эталонные решения, метод
изготовленных решений, контроль сеточной сходимости, законов сохранения и
независимые реализации. Валидация требует экспериментальных или надёжных
наблюдательных данных. Поэтому верификация логически должна предшествовать
валидации.

\subsection{Сопоставление с экспериментом}

Сравнивать желательно не отдельную точку, а весь диапазон предполагаемого
применения модели. Помимо максимальной ошибки используют, например,
\[
    \mathrm{RMSE}
    =\sqrt{\frac1N\sum_{i=1}^{N}r_i^2}.
\]
Важно анализировать не только величину, но и \textbf{структуру остатков}: их
систематический знак, тренд по входным параметрам или времени часто указывает на
пропущенный физический эффект.

Если параметры $\theta$ подбирались по данным $D_{\rm cal}$, прогнозирующую
способность модели желательно проверять на независимых данных $D_{\rm val}$.
Иначе хорошее совпадение может быть лишь следствием настройки параметров под
конкретный набор наблюдений.

\subsection{Статистическая проверка адекватности регрессионной модели}

В планировании эксперимента полезно отделять разброс, вызванный случайной
\textbf{ошибкой воспроизводимости}, от систематического расхождения модели и
средних экспериментальных значений.

Пусть в $N$ различных точках проведено по $n_i$ повторных опытов,
$y_{ij}$ --- результаты, а
\[
    \overline y_i=\frac1{n_i}\sum_{j=1}^{n_i}y_{ij},
    \qquad
    \widehat y_i=y^{\rm mod}(x_i,\widehat\theta).
\]
Тогда сумма квадратов ошибки воспроизводимости
\[
    SS_{\rm err}
    =\sum_{i=1}^{N}\sum_{j=1}^{n_i}
      (y_{ij}-\overline y_i)^2,
    \qquad
    f_{\rm err}=\sum_{i=1}^{N}(n_i-1),
\]
а сумма квадратов неадекватности
\[
    SS_{\rm ad}
    =\sum_{i=1}^{N}n_i(\overline y_i-\widehat y_i)^2,
    \qquad
    f_{\rm ad}=N-p,
\]
где $p$ --- число оценённых коэффициентов модели. При стандартных предположениях
о независимых нормальных ошибках с одинаковой дисперсией сравнивают
\[
    \boxed{
    F=\frac{SS_{\rm ad}/f_{\rm ad}}
            {SS_{\rm err}/f_{\rm err}}
    }
\]
с критическим значением распределения Фишера. Если $F$ не превышает критического
уровня, статистических оснований отвергать гипотезу об адекватности модели нет.
Если превышает, необходимо пересмотреть структуру модели, диапазон её применения
или качество исходных данных.

В простейшем варианте повторные опыты проводят в центральной точке плана:
расчётное значение должно согласовываться с экспериментальным средним в пределах
его доверительного интервала. Такой подход прямо связывает понятие адекватности с
ошибкой воспроизводимости эксперимента.

\subsection{Калибровка, идентифицируемость и чувствительность}

Параметры часто определяют из задачи
\[
    \widehat\theta
    =\arg\min_{\theta}
    \sum_i w_i
    \bigl(y_i^{\rm exp}-y_i^{\rm mod}(\theta)\bigr)^2.
\]
Это \textbf{калибровка}, а не самостоятельное доказательство адекватности.
Модель может хорошо описывать калибровочные данные при существенно разных
значениях параметров.

\textbf{Идентифицируемость} означает возможность однозначно восстановить
интересующие параметры по наблюдаемым величинам. Локальную чувствительность
выхода $y$ к параметрам характеризует матрица
\[
    S_{ij}=\frac{\partial y_i}{\partial\theta_j}.
\]
Малочувствительные параметры трудно надёжно определять по эксперименту, а почти
линейно зависимые столбцы $S$ указывают на практическую неидентифицируемость.

\subsection{Неопределённость и область применимости}

В итоговом сравнении следует учитывать по меньшей мере четыре источника
неопределённости:
\begin{itemize}
    \item погрешности экспериментальных данных;
    \item неопределённость параметров;
    \item погрешность численного метода и вычислений;
    \item структурную погрешность самой математической модели.
\end{itemize}

Успешная проверка в конечном наборе режимов не превращает модель в универсально
верную. Корректный вывод формулируется примерно так: \emph{модель согласуется с
данными с заданной точностью для указанных характеристик и в исследованном
диапазоне параметров}.

\subsection{Что важно сказать на экзамене}

Нужно определить адекватность относительно \textbf{цели, точности и области
применимости}, затем обязательно различить
\[
\boxed{\text{верификация}\;\neq\;\text{валидация}\;\neq\;\text{калибровка}}.
\]
Для экспериментальной регрессионной модели полезно объяснить разделение ошибки
воспроизводимости и ошибки неадекватности и критерий Фишера. В конце следует
упомянуть независимые проверочные данные, идентифицируемость параметров и то, что
область адекватности всегда ограничена проверенными режимами.

\newpage
